\section{Preliminary of Automata and Krohn-Rhodes Decomposition Theorem}\label{appsec:prelim_automata}

In this section we recap the basic notations and definitions for automata theory and Krohn-Rhodes Decomposition Theorem~\citep{krohn1965algebraic}, following the notation and presentation of \citet{maler2010krohn}. 

\begin{definition}[Automaton]\label{defi:automaton}
	A deterministic automaton is triple $\gA = (\Sigma,Q,\delta)$ where
$\Sigma$ is a finite set of symbols called the input alphabet, $Q$ is a finite set of states and
$\delta : Q \times \Sigma \to Q$ is the transition function.
\end{definition}

The transition function can be lifted naturally to input sequences, by letting $\delta(q, w\sigma) \triangleq \delta(\delta(q, w), \sigma)$ for all $w\in \Sigma^*$ recursively.

An automaton can be made an acceptor by choosing an initial state $q_0 \in Q$ and a set of accepting states $F \subseteq Q$. As such it accepts/recognizes a set of sequences, also known as a language, defined as $\gL(A,q_0,F) = \{  w\in\Sigma^* : \delta(q_0, w) \in F\}$. Kleene’s Theorem states that the class of languages recognizable by finite automata coincides with the regular
languages.


\begin{definition}[Automaton Homomorphism]
	A surjection $\phi : Q \to Q_0$
is an automaton homomorphism from $\gA = (\Sigma, Q, \delta)$ to $\gA_0 = (\Sigma, Q_0, \delta_0)$ if for every $q \in Q, \sigma \in \Sigma,\phi(\delta(q, \sigma)) = \delta_0(\phi(q), \sigma)$.
In such a case we say that $\gA_0$
is homomorphic to $\gA$ and denote it by $\gA_0 \le_\phi \gA$. When $\phi$
is a bijection, $\gA$ and $\gA_0$ are said to be isomorphic.
\end{definition} 

The conceptual significance of Automaton Homomorphism is that, if we can simulate any $\gA$ and $\gA_0 \le_\phi \gA$, we can `almost' simulate $\gA_0$ as well, in the sense of following lemma:

\begin{lemma}\label{lem:automata_homo}
	For any two automata $\gA= (\Sigma, Q, \delta),\gA_0 = (\Sigma, Q_0, \delta_0)$ satisfying that $\gA_0 \le_\phi \gA$ for some function $\phi$,  for any $F_0\subseteq Q$, $q_0\in Q$, $\phi(q) = q_0$, it holds that $ \gL(\gA_0,q_0,F_0) = \gL(\gA,q,\phi^{-1}(F_0))$.
\end{lemma}

\begin{proof}[Proof of \Cref{lem:automata_homo}]
We claim for any $w\in\Sigma^*$, it holds that $\phi(\delta(q,w)) = \delta(\phi(q),w)$. This claim holds by definition of automaton homomorphism for all $|w|\le 1$. suppose the claim already holds for all $w$ no longer than $n$ for some $n$, for any $w'=w\sigma$ with $|w|=n$ and $\sigma\in\Sigma$, it holds that $\phi(\delta(q,w')) = \phi(\delta(\delta(q,w),\sigma)) = \delta(\phi(\delta(q,w)),\sigma)=\delta(\delta(\phi(q),w),\sigma) = \delta(\phi(q),w')$. 	Therefore $\delta_0(q_0,w)\in F_0 \iff \delta_0(\phi(q),w)\in F_0 \iff \phi(\delta(q,w))\in F_0 \iff \delta(q,w)\in \phi^{-1}(F_0)$. Thus we conclude that 
$\gL(\gA_0,q_0,F_0) = \{w\in\Sigma^*\mid \delta(q_0,w)\in F_0\} = \{w\in\Sigma^*\mid \delta(q,w)\in \phi^{-1}(F_0)\} = \gL(\gA,q,\phi^{-1}(F_0))$.
\end{proof}


\begin{definition}[Semigroups, Monoids and Groups]
	A Semigroup is a pair $(S, \cdot)$ where
$S$ is a set and $\cdot$ is a binary associative operation (“multiplication”) from $S\times S$ to $S$. A
Monoid $(S, \cdot, 1)$ is a semigroup admitting an identity element $1$ such that $s\cdot1 = 1\cdots = s$
for every $s \in S$. A group is a monoid such that for every $s \in S$ there exists an element
$s^{-1} \in S$ (an inverse) such that $s \cdot s^{-1} = 1$.
\end{definition}

\begin{definition}[Semigroup Homomorphisms]
	A surjective function $\phi : S \to S_0$
is a semigroup homomorphism from $(S, \cdot)$ to $(S_0, \ast)$ if for every $s_1,s_2 \in S, 
\phi(s_1 \cdot s_2) = \phi(s_1) \ast \phi(s_2)$.
In such a case we say that $S_0$ is homomorphic to $S$ and denote it by $S_0 \le_\phi S$. Two mutually homomorphic semigroups are said to be isomorphic.
\end{definition}


\begin{definition}[Transformation Semigroup]
	The transformation semigroup of an automata $\gA = (\Sigma,Q,\delta)$ is the semigroup generated by $\{\delta(\cdot, \sigma):Q\to Q\mid \sigma \in \Sigma \}$.
\end{definition}




\subsection{The Krohn-Rhodes Decomposition Theorem}\label{appsec:krhon-rhodes theorem}

Below we give the definition of the cascade product of two automata, which is a central concept used in Krohn-Rhodes Decomposition Theorem for automata.

\begin{definition}[Cascade Product]
Let $\gB_1 = (\Sigma, Q_1, \delta_1)$ and 
$\gB_2 = (Q_1 \times \Sigma, Q_2, \delta_2)$ be two automata. The cascade product $\gB_1 \circ \gB_2$ is the automaton $C = (\Sigma, Q_1\times Q_2, \overline\delta)$ where
$$\overline\delta((q_1, q_2), \sigma) \triangleq (\delta_1(q_1, \sigma), \delta_2(q_2,(q_1, \sigma))).$$
The cascade product of more than two automata is defined as
$\gB_1 \circ \gB_2 \circ \cdots\circ \gB_k = (\cdots((\gB_1 \circ \gB_2) \circ B_3\cdots) \circ \gB_k$.
\end{definition}

\begin{definition}[Permutation-Reset Automata]
A automaton $\gA =(\Sigma, Q, \delta)$ is a permutation-reset automaton if for every letter $\sigma \in \Sigma$, $\sigma$ is either a
permutation or reset. If the only
permutations are identities, we call it a reset automaton.
\end{definition}
\begin{theorem}[Krohn-Rhodes; cf. \citet{maler2010krohn}]\label{thm:krohn_rhodes}
	For every automaton A there exists a cascade $\gC = \gB_1 \circ \gB_2 \circ \cdot \cdot \cdot \circ \gB_k$ such that:
	\begin{enumerate}
		\item Each $\gB_i$ is a permutation-reset automaton;
		\item There is a homomorphism $\phi$ from $\gC$ to $\gA$;
		\item Any permutation group in some $\gB_i$
is homomorphic to a subgroup of the transformation semigroup of $\gA$.

	\end{enumerate}
The pair $(\gC, \phi)$ is called a cascaded decomposition of $\gA$.
\end{theorem}

 
\subsection{Counter-free Automata}\label{appsec:counter_free_automata}

Next we introduce a key concept used in the proof of \Cref{thm:iterated_addition_in_AC0} (and thus \Cref{thm:AC0_upper_bound}) -- Counter-free Automaton.
\begin{definition}[Counter-free Automaton, \citep{mcnaughton1971counter}]
	An automaton is counter-free if no word $w\in\Sigma^*$ induces a permutation other than identity on any subset of $Q$. 
\end{definition}

 A subclass of the regular languages is the class of star-free sets defined as:
\begin{definition}[Star-Free regular languages]\label{defi:star_free_language}
	 The class of star-free regular languages over $\Sigma$ is the
smallest class containing $\Sigma^\ast$ and the sets of the form $\{  \sigma\}$ where $\sigma \in \sigma \cup \{ \eps\}$, which is
closed under finitely many applications of concatenation and Boolean operations including union, intersection, and complementation.
\end{definition}

It is well-known that languages recognized by counter-free automata have the following equivalent characterizations.

\begin{theorem}[\citet{mcnaughton1971counter}]\label{thm:counter_free_language_three_characterization}
Suppose $L$ is a regular language not containing the empty string. Then the following are equivalent:
 \begin{enumerate}
 	\item $L$ is star-free;
 	\item $L$ is accepted by a counter-free automata.
 	\item $L$ is non-counting, \emph{i.e.}, there is an $n\in\mathbb{N}$ so that for all $x$, $y$, and $z$ and all $m \ge n$, $xy^mz \in  L  \iff xy^{m+1}z\in L$. 
 \end{enumerate}
\end{theorem}

Counter-free property of an automaton can also be characterized via its transformation semigroup by \Cref{lem:aperiodic}, whose proof is straightforward and skipped. 

\begin{lemma}\label{lem:aperiodic}
An automaton is counter-free if and only if the transformation semigroup of the automaton is group-free, \emph{i.e.}, it has no non-trivial subgroups. A semigroup $(S,\cdot)$ is group-free if and only if it is \emph{aperiodic}, \emph{i.e.}, for all $s\in S$, there exists $k\in\mathbb{N}$, $s^k = s^{k+1}$.
\end{lemma}


Thus \Cref{cor:krohm_rhodes_aperiodic} holds as a corollary of \Cref{thm:krohn_rhodes}.

\begin{theorem}[Corollary of \Cref{thm:krohn_rhodes}]\label{cor:krohm_rhodes_aperiodic}
		For every counter-free automaton $\gA$ there exists a cascade $\gC = \gB_1 \circ \gB_2 \circ \cdot \cdot \cdot \circ \gB_k$ such that each $\gB_i$ is a reset automaton and there is a homomorphism $\phi$ from $\gC$ to $\gA$.
\end{theorem}


Using \Cref{cor:krohm_rhodes_aperiodic} the following theorem connects the counter-free automata to constant-depth poly-size circuits with unbounded fan-in. The high-level proof idea is that any reset automaton can be simulated using constantly many depth and any counter-free automaton can be decomposed into the cascade product of a finite number of reset automaton. 

\begin{theorem}\label{thm:counter_free_in_AC0}[Theorem 2.6, \citet{chandra1983unbounded}]
	Suppose $\gA=(\Sigma,Q,\delta)$ is an counter-free automaton. Then there is a circuit of size $O(n^3)$ with unbounded fan-in 
and constant depth that simulates $\delta(q,w)$ for any $q\in Q$ and $w\in\Sigma^*$ satisfying $|w|=n$, where $O(\cdot)$ hides constants depending on the automaton. 
\end{theorem}
